% !TeX root = ../../Book1.tex
\OptionalSection{绝对连续曲线、度量导数与长度}{b1:sec:1605}

在欧氏空间中，曲线的导数同时记录速度的大小与方向。一般度量空间没有向量差，方向未必能够定义，但相邻时刻的位置仍有距离。用距离除以时间差，可以保留速度的大小；要使它与长度的积分相联系，合适的条件是绝对连续性。

\ProseParagraphBreak

本节的参数区间固定为 \([a,b]\)，其中 \(a<b\)。空间只假设为度量空间，不要求线性、完备或局部紧。上一节已经建立实值绝对连续函数的积分表示；这里把一般曲线转换成可数个实值距离函数，使测度论微积分能够用于没有坐标的空间。

% 度量空间绝对连续性定义对照 Fremlin2016Measure 作者官方 chap26.tex 的 264Yp。
% 度量导数对照 Rossi--Mielke--Savare (2008)，A metric approach to a class of doubly
% nonlinear evolution equations and applications，作者公开版 §2.1，Proposition 2.2 及证明，pp.10--12：
% https://riccarda-rossi.unibs.it/pubblicazioni/Published%20papers/pubblicazioni-files/rossi-mielke-savare.pdf
% 正式发表信息：Ann. Scuola Norm. Sup. Pisa Cl. Sci. (5), 7(1), 97--169。
% 该文处理更广的非对称距离；本节只证对称度量，用可数距离测试和长度函数独立组织证明。

\subsection{绝对连续曲线}
\label{b1:subsec:160501}

\begin{definition}
\label{b1:def:absolutely-continuous-curve}
设 \((X,d)\) 为度量空间，\(\gamma:[a,b]\to X\)。若对每个 \(\varepsilon>0\)，存在 \(\delta>0\)，使任意有限族内点互不相交的区间 \([s_i,t_i]\subseteq[a,b]\) 满足
\[
\sum_i(t_i-s_i)<\delta
\quad\Longrightarrow\quad
\sum_i d\bigl(\gamma(s_i),\gamma(t_i)\bigr)<\varepsilon,
\]
则称 \(\gamma\) 为\term[juedui-lianxu-quxian]{绝对连续曲线}[absolutely continuous curve]，记作 \(\gamma\in\operatorname{AC}([a,b];X)\)。当 \(X=\mathbb R\) 时，这一定义与上一节的绝对连续实值函数一致。
\end{definition}

这里只取一个区间，便可看出绝对连续曲线一致连续；它的像是紧的，因而是可分度量空间。Lipschitz 曲线满足定义，因为所有距离之和不超过 Lipschitz 常数乘以区间长度之和。绝对连续性允许不同位置具有不同的速度控制，但要求小的总时间长度只能产生小的总位移。

\begin{lemma}
\label{b1:lem:ac-variation}
设 \(\gamma\in\operatorname{AC}([a,b];X)\)，定义
\[
V(t)=\sup\Bigl\{\sum_{i=1}^n d\bigl(\gamma(t_i),\gamma(t_{i-1})\bigr):
a=t_0<\cdots<t_n=t\Bigr\},
\quad V(a)=0.
\]
则 \(V\) 有限、单调不减且绝对连续。对 \(a\leq s<t\leq b\)，有
\[
d\bigl(\gamma(s),\gamma(t)\bigr)\leq V(t)-V(s),
\]
而右端正是同一个分割上确界在子区间 \([s,t]\) 上的值。
\end{lemma}
\begin{proof}
在一个分割中插入中间点，由三角不等式，距离和不会减小；反过来，可以拼接两个子区间的分割。这证明分割上确界对子区间可加，允许暂取无穷值。

对 \(\varepsilon=1\) 取定义中的 \(\delta\)，把 \([a,b]\) 分成有限个长度小于 \(\delta\) 的区间。在其中任一区间内，任意分割的各小段总长度小于 \(\delta\)，所以距离和小于 \(1\)。把一个任意分割加上这些预定端点，便得到总上界，故 \(V(b)<\infty\)。单调性、子区间公式及位移估计随之成立。

最后，给定 \(\varepsilon>0\)，对绝对连续性取使距离和小于 \(\varepsilon/2\) 的 \(\delta\)。若一族不交区间的总长度小于 \(\delta\)，在每个区间内部选分割来逼近其分割上确界。所有分割小段仍组成总长度小于 \(\delta\) 的不交族。因此，取这些有限个上确界后，
\[
\sum_i\bigl(V(t_i)-V(s_i)\bigr)\leq\varepsilon/2<\varepsilon.
\]
这正是 \(V\) 的绝对连续性。
\end{proof}

由\cref{b1:lem:ac-variation}与\cref{b1:thm:ac-fundamental-calculus}立即得到一个等价而实用的条件：\(\gamma\) 绝对连续，当且仅当存在非负 \(g\in L^1([a,b])\)，使
\[
d\bigl(\gamma(s),\gamma(t)\bigr)
\leq\int_s^t g(r)\,\mathrm dr
\quad(a\leq s<t\leq b).
\]
正向由\cref{b1:lem:ac-variation}知 \(V\) 绝对连续，故微积分基本定理给出
\(V(t)-V(s)=\int_s^tV'(r)\,\mathrm dr\)，并且 \(V'\geq0\) 几乎处处；取 \(g=V'\) 即可。
反向对不交区间求和，再用积分绝对连续性。于是绝对连续曲线确有一个可积的速度上界。下一小节将证明，这些上界中有一个几乎处处最小的选择。

\subsection{度量导数}
\label{b1:subsec:160502}

\begin{definition}
\label{b1:def:metric-derivative}
设 \(\gamma:[a,b]\to X\) 为曲线。若在 \(t\in(a,b)\) 处极限
\[
|\gamma'|_d(t)
\coloneqq\lim_{h\to0}\frac{d\bigl(\gamma(t+h),\gamma(t)\bigr)}{|h|}
\]
存在且有限，则称其为 \(\gamma\) 在 \(t\) 处的\term[duliang-daoshu]{度量导数}[metric derivative]，也称该处的度量速度。这里 \(|\gamma'|_d\) 是一个整体记号，并不预设 \(\gamma'\) 作为向量存在。
\symidx{\(\lvert\gamma'\rvert_d(t)\)：曲线的度量导数}
\end{definition}

度量导数及其最小控制性质，可以与 Rossi、Mielke、Savar\'e 的《A metric approach to a class of doubly nonlinear evolution equations and applications》\cite[Section 2.1, Proposition 2.2]{Rossi2008Metric}对照阅读。该文还允许非对称距离，服务于演化方程的度量表述；本节只讨论通常的对称度量，并将证明组织为距离函数与长度函数之间的比较。

\begin{theorem}[度量导数的存在性]
\label{b1:thm:metric-derivative}
若 \(\gamma\in\operatorname{AC}([a,b];X)\)，则其度量导数几乎处处存在。在例外零集上把它赋值为零后，它是可积函数，并且
\[
d\bigl(\gamma(s),\gamma(t)\bigr)
\leq\int_s^t|\gamma'|_d(r)\,\mathrm dr
\quad(a\leq s<t\leq b).
\]
若非负 \(g\in L^1([a,b])\) 也满足这个位移控制不等式，则 \(|\gamma'|_d\leq g\) 几乎处处。
\end{theorem}
\begin{proof}
在紧像 \(\gamma([a,b])\) 中取可数稠密集 \(\{y_j:j\geq1\}\)，令
\[
f_j(t)=d\bigl(\gamma(t),y_j\bigr).
\]
三角不等式给出 \(|f_j(t)-f_j(s)|\leq d\bigl(\gamma(t),\gamma(s)\bigr)\)，因此每个 \(f_j\) 都是实值绝对连续函数。它们还完整记录像中两点的距离：
\[
d(x,y)=\sup_j|d(x,y_j)-d(y,y_j)|
\quad(x,y\in\gamma([a,b])).
\]
其中“\(\geq\)”来自三角不等式；为得反向不等式，只需选 \(y_j\to x\)。

在全部 \(f_j\) 及引理中的 \(V\) 都可微的共同满测度集上，定义
\[
v(t)=\sup_j|f_j'(t)|,
\]
在其补集上赋值为零。导数可由可测差商的极限得到，因此 \(v\) 可测。由 \(|f_j(t)-f_j(s)|\leq|V(t)-V(s)|\)，有 \(0\leq v\leq V'\) 几乎处处，故 \(v\in L^1\)。实值微积分基本定理又给出
\[
|f_j(t)-f_j(s)|\leq\int_s^t v(r)\,\mathrm dr.
\]
对 \(j\) 取上确界，得到 \(d\bigl(\gamma(s),\gamma(t)\bigr)\leq\int_s^tv\)。在 \([s,t]\) 的任意分割上求和并取上确界，进一步得到
\[
V(t)-V(s)\leq\int_s^t v(r)\,\mathrm dr.
\]
在 \(V\) 可微且 \(v\) 为 Lebesgue 点之处取差商，可知 \(V'\leq v\) 几乎处处。因此
\[
v=V'\quad\text{几乎处处}.
\]

现在取一个满足上述等式、且全部 \(f_j\) 和 \(V\) 可微的内部点 \(t\)。对每个 \(j\)，距离商不小于 \(|f_j(t+h)-f_j(t)|/|h|\)，所以
\[
\liminf_{h\to0}\frac{d\bigl(\gamma(t+h),\gamma(t)\bigr)}{|h|}
\geq\sup_j|f_j'(t)|=v(t).
\]
另一方面，由\cref{b1:lem:ac-variation}，同一个距离商不超过 \(|V(t+h)-V(t)|/|h|\)，故其上极限不超过 \(V'(t)=v(t)\)。这证明极限存在，且 \(|\gamma'|_d=V'=v\) 几乎处处，并给出可积性及位移估计。

若 \(g\) 是另一个可积控制函数，则在度量导数存在且 \(g\) 为 Lebesgue 点的地方，对 \(\gamma(t+h)\) 与 \(\gamma(t)\) 使用控制不等式，再除以 \(|h|\) 并令 \(h\to0\)，便得到 \(|\gamma'|_d(t)\leq g(t)\)。
\end{proof}

这个证明的可数性来自曲线的紧像，而不是来自环境空间的可分性。它也没有把可数个实导数组合成一个向量导数，只是用这些导数确定了一个非负速度。

\ProseParagraphBreak

这种区别在无限维空间中确实会发生。取 \(X=L^1(0,1)\)，令 \(\gamma(t)=\mathbf1_{[0,t]}\)。则
\[
\|\gamma(t)-\gamma(s)\|_1=|t-s|,
\]
所以曲线是等距的，\(|\gamma'|_d(t)=1\) 在所有内部点成立。然而，对充分小的 \(h>0\)，正向差商为 \(q_h=h^{-1}\mathbf1_{(t,t+h]}\)，并且
\[
\|q_h-q_{h/2}\|_1
=\int_t^{t+h/2}\frac1h\,\mathrm dr
+\int_{t+h/2}^{t+h}\frac1h\,\mathrm dr=1.
\]
差商不是 Cauchy 族，故在任何内部点都不存在 \(L^1\) 范数意义下的向量导数。一般 Banach 空间中的绝对连续曲线不能直接套用有限维 Rademacher 定理；度量导数的存在性却不受这个障碍影响。

\subsection{长度公式}
\label{b1:subsec:160503}

\begin{definition}
\label{b1:def:curve-length}
给定连续曲线 \(\gamma:[a,b]\to X\)，称
\[
\ell(\gamma)
\coloneqq\sup\Bigl\{\sum_{i=1}^n d\bigl(\gamma(t_i),\gamma(t_{i-1})\bigr):
a=t_0<\cdots<t_n=b\Bigr\}
\]
为曲线的\term[quxian-changdu]{长度}[length of a curve]。若 \(\ell(\gamma)<\infty\)，则称其为\term[keqiuchang-quxian]{可求长曲线}[rectifiable curve]。
\symidx{\(\ell(\gamma)\)：参数化曲线的长度}
\end{definition}

长度计入曲线的行进过程，不只是像集的大小。绝对连续曲线的长度就是先前的 \(V(b)\)，而 \(V(t)\) 记录截至时刻 \(t\) 的累积长度。

\begin{theorem}[绝对连续曲线的长度公式]
\label{b1:thm:metric-curve-length}
设 \(\gamma\in\operatorname{AC}([a,b];X)\)。则对任意 \([s,t]\subseteq[a,b]\)，
\[
\ell\bigl(\gamma|_{[s,t]}\bigr)
=\int_s^t|\gamma'|_d(r)\,\mathrm dr.
\]
若 \(X=\mathbb R^m\)，则 \(\gamma'\) 几乎处处存在，且
\[
|\gamma'|_d=\|\gamma'\|
\quad\text{几乎处处},
\quad
\ell(\gamma)=\int_a^b\|\gamma'(t)\|\,\mathrm dt.
\]
\end{theorem}
\begin{proof}
度量导数定理的证明已经得到 \(|\gamma'|_d=V'\) 几乎处处。由 \(V\) 的绝对连续性及子区间可加性，
\[
\ell\bigl(\gamma|_{[s,t]}\bigr)
=V(t)-V(s)=\int_s^tV'(r)\,\mathrm dr,
\]
这就是长度公式。若像在 \(\mathbb R^m\) 中，每个坐标函数都是实值绝对连续函数，因而可微的共同满测度集只需除去有限个零集。在此集合上，向量差商收敛到坐标导数组成的 \(\gamma'\)，取范数便得 \(|\gamma'|_d=\|\gamma'\|\)。
\end{proof}

绝对连续性不能由“连续且可求长”替换。Cantor 函数作为 \([0,1]\) 中的一条单调连续曲线，长度为 \(1\)，度量速度却几乎处处为零。它把一部分行进过程集中在零测参数集上，导数的 Lebesgue 积分不能恢复这部分长度。

\begin{proposition}
\label{b1:prop:arc-length-parameterization}
设 \(\gamma\in\operatorname{AC}([a,b];X)\)，且 \(L=\ell(\gamma)>0\)。则存在 \(1\)-Lipschitz 曲线 \(\eta:[0,L]\to X\)，满足
\[
\gamma(t)=\eta\bigl(V(t)\bigr),
\quad
\ell\bigl(\eta|_{[u,v]}\bigr)=v-u
\quad(0\leq u<v\leq L).
\]
因而 \(|\eta'|_d=1\) 几乎处处。这称为曲线的\term[huchang-canshuhua]{弧长参数化}[arc-length parameterization]。
\end{proposition}
\begin{proof}
函数 \(V\) 连续、单调不减，从 \([a,b]\) 映满 \([0,L]\)。定义 \(\eta(s)=\gamma(t)\)，其中 \(V(t)=s\)。若两个参数有相同的 \(V\) 值，则它们的像点距离不超过 \(V\) 的增量零，所以定义与选择无关。同一个位移估计给出 \(d\bigl(\eta(u),\eta(v)\bigr)\leq v-u\)，故 \(\eta\) 是 \(1\)-Lipschitz 的，特别有 \(\ell(\eta)\leq L\)。

反过来，把 \(\gamma\) 的任意分割经 \(V\) 映到 \([0,L]\)，删去重复参数，距离和保持不变。因此 \(\ell(\gamma)\leq\ell(\eta)\)，遂有 \(\ell(\eta)=L\)。长度对相邻子区间可加，而每个子区间的长度都不超过其参数长度，于是 \([0,u]\)、\([u,v]\)、\([v,L]\) 三段必须分别取到这个上界。故 \(\ell\bigl(\eta|_{[u,v]}\bigr)=v-u\)。由长度公式，\(\int_u^v|\eta'|_d=v-u\) 对所有子区间成立，再用 Lebesgue 微分即得 \(|\eta'|_d=1\) 几乎处处。
\end{proof}

若 \(L=0\)，曲线由位移估计必为常值。对 \(L>0\)，弧长参数化与上一章的 Lipschitz 测度估计还给出
\[
\mathcal H^1\bigl(\gamma([a,b])\bigr)
=\mathcal H^1\bigl(\eta([0,L])\bigr)
\leq\mathcal H^1([0,L])=L.
\]
这里用到了\cref{b1:prop:hausdorff-lipschitz}及\cref{b1:thm:hausdorff-lebesgue}。不等式可以严格：曲线先从 \(0\) 匀速走到 \(1\)，再原路返回 \(0\)，其像为 \([0,1]\)，一维测度为 \(1\)，但长度为 \(2\)。像集测度不计重复经过，长度则把它计入；下一章的面积公式将以重数处理这种区别。
